\chapter*{Введение}
\addcontentsline{toc}{chapter}{Введение}
\label{cha:intro}
%\Introduction

Паллетизация остаётся ключевым технологическим узлом современной складской логистики: именно от плотности и устойчивости укладки зависит стоимость хранения, расход упаковочных материалов и объём перевозимого груза. С ростом ассортимента e-commerce-складов задача становится всё более вариативной, а ручная разработка правил паллетизации — трудоёмкой и часто приводит к недогрузам до 50 \%.

\textbf{Цель работы} — разработка и программная реализация семейства алгоритмов трёхмерной укладки коробок на паллету, обеспечивающих максимальный коэффициент заполнения при производительном времени расчёта.

В первом разделе работы содержится постановка задачи  и подзадач, на которые она была разбита.

Во втором разделе работы приведен обзор научно - технической литературы по теме исследований.  

Третий раздел формализует задачу и описывает исходные данные, используемые для построения алгоритмов.
 
В четвертом разделе приведены описания разработанных алгоритмов. 

В пятом разделе приведено сравнение разработанных алгоритмов. 